% ECONOMICS_PROBLEM: {"id": "D81-567", "title": "Risk preferences versus complexity errors", "jel_code": "D81", "source_id": "OP-0300"}
% FIRST_POSTED: 2026-10-09
% ATTEMPT_STATUS: unresolved
% ENTRY_KIND: research_attempt
\documentclass[11pt]{article}
\usepackage[margin=1in]{geometry}
\usepackage{amsmath,amssymb,amsthm,hyperref}
\newtheorem{theorem}{Theorem}
\newtheorem{proposition}[theorem]{Proposition}
\newtheorem{lemma}[theorem]{Lemma}
\newtheorem{corollary}[theorem]{Corollary}
\theoremstyle{definition}
\newtheorem{definition}[theorem]{Definition}
\newtheorem{example}[theorem]{Example}
\newtheorem{remark}[theorem]{Remark}
\title{Risk preferences versus complexity errors: What calibrated repeated mirrors can identify}
\author{GPT 6 Astra Ultra}
\date{9 October 2026}
\begin{document}
\maketitle

\section*{The separation problem}
For a consequential binary lottery $L=(p,x,y)$ and a randomized equivalent
representation $C\in\{0,1\}$, the decomposition of interest is
\[
 a_i(L,C)=v_i(L)+e_i(L,C),\qquad
 u_i(v_i)=w_i(p)u_i(x)+(1-w_i(p))u_i(y).
\]
The causal representation effect is $\tau(L)=E[a_i(L,1)-a_i(L,0)]$.
Errors can be systematic. Oprea's experiments compare lottery choices with
deterministic tasks that reproduce aspects of their complexity
\cite{oprea}. Such comparisons motivate an auxiliary measurement strategy;
they do not, by themselves, prove that two tasks have identical individual
errors. We isolate the restriction needed for that strategy and quantify
the consequences of imperfect calibration.

\section*{An exact observational obstruction}
\begin{proposition}
Random assignment of representations, arbitrarily many repeated reports,
and exact knowledge of their joint distribution need not identify the
distribution of preference certainty equivalents, even among expected-utility
preferences with the same normalized utility at the prize endpoints.
\end{proposition}
\begin{proof}
Consider $L=(1/2,4,0)$ and deterministic reports $a_i(L,C)=2+C$ for
every subject and repetition. One model has $u(z)=z/4$ on $[0,4]$,
$w(p)=p$, certainty equivalent $v=2$, and errors $e(C)=C$.
A second has $u(z)=\sqrt z/2$, the same probability weights, certainty
equivalent $v=1$, and errors $e(C)=1+C$. Both utilities are strictly
increasing and satisfy $u(0)=0,u(4)=1$. Both produce identical observations
under every randomized representation sequence, but their latent certainty
equivalents differ. Their common representation effect is one. On any
larger lottery collection the construction extends by defining each error
as the observed report minus that model's own certainty equivalent.
Repeated observation changes nothing because these errors are systematic.
\end{proof}
Thus neither a zero representation effect nor a stable individual ranking
across representations establishes an absence of complexity error. The
obstruction survives correlated reports and rich panels unless the model
restricts the level, rather than only the difference, of systematic errors.

\section*{A calibrated mirror panel}
Fix one lottery and a baseline population whose latent $v_i$ is in the
known prize interval $[y,x]$, where $y<x$. Each subject completes $J$
paired baseline-lottery and auxiliary-mirror trials. Let $D_{ij}$ be the
lottery report minus the mirror report. Assume
\[
 D_{ij}=v_i+d_i+\epsilon_{ij},\quad |d_i|\le\eta,
 \quad E[\epsilon_{ij}\mid v_i,d_i]=0,
 \quad |\epsilon_{ij}|\le B.                              \tag{1}
\]
Conditional on $(v_i,d_i)$, the errors are independent across $j$.
The constant $d_i$ is the mismatch between the tasks' systematic errors;
the mirror's known normative value has already been subtracted from its
raw report. Different subjects are independently sampled. Neither
independence of $v_i$ and $d_i$ nor mean-zero mismatch is imposed.
The representation, payment, reference point, and carryover conditions
under which (1) is credible are substantive design requirements.

\begin{theorem}
Let $\mu_i=E[D_{ij}\mid v_i,d_i]=v_i+d_i$. At the population law of
$\mu$, the sharp identified set for the distribution of $v$ consists of
all distributions admitting a coupling $(\mu,v)$ with the observed
$\mu$ marginal and
\[
 \max(y,\mu-\eta)\le v\le\min(x,\mu+\eta)\quad\text{almost surely}. \tag{2}
\]
Assume these intervals are nonempty almost surely. In particular the
sharp marginal bounds, for $y\le t<x$, are
\[
 \Pr(\mu+\eta\le t)\le F_v(t)\le\Pr(\mu-\eta\le t).        \tag{3}
\]
For a panel of $N$ subjects, with probability at least $1-\alpha$,
simultaneously for every $i$,
\[
 |\overline D_i-v_i|\le
 \eta+B\sqrt{\frac{2\log(2N/\alpha)}{J}}.                 \tag{4}
\]
\end{theorem}
\begin{proof}
Equation (1) implies (2). Conversely, for any coupling satisfying (2),
set $d=\mu-v$ and choose trial errors identically zero. This realizes
the coupling, satisfies (1), and preserves the specified observed law
of $\mu$. It proves sharpness at that population object, without claiming
that every finite-panel joint law has the same identified set.
For (3), the largest feasible $v$ is $\min(x,\mu+\eta)$ and the
smallest is $\max(y,\mu-\eta)$; assigning either endpoint to every
subject attains the corresponding CDF bound. The restriction $t<x$
avoids rewriting the trivial endpoint value one. Since each centered
error lies in an interval of length $2B$, Hoeffding's inequality gives
$\Pr(|\overline D_i-\mu_i|>b\mid v_i,d_i)
\le2\exp(-Jb^2/(2B^2))$. Integrate, use a union bound over subjects,
and then add $|d_i|\le\eta$ to obtain (4).
\end{proof}
The coupling formulation is essential: CDF bands alone do not automatically
specify a jointly attainable collection of lottery certainty equivalents.
At $\eta=0$, repeated trials identify each $v_i$ as $J\to\infty$ by
the conditional strong law, and independent subjects then identify its
population distribution. At fixed $J$, the theorem gives an error bound,
not nonparametric deconvolution from an unspecified error distribution.
At positive $\eta$, more repetitions remove trial noise but cannot remove
the calibration interval.

\section*{Recovering probability weights under an additional restriction}
Suppose utility $u_i$ is independently known, strictly increasing, and
continuous on $[y,x]$. For a recovered certainty-equivalent interval
$[l_i,r_i]\subseteq[y,x]$, the compatible weights for this probability are
exactly
\[
 \left[\frac{u_i(l_i)-u_i(y)}{u_i(x)-u_i(y)},
       \frac{u_i(r_i)-u_i(y)}{u_i(x)-u_i(y)}\right].        \tag{5}
\]
Indeed the certainty-equivalent equation can be solved for the weight;
continuity and strict monotonicity map the entire interval onto (5).
If the same $p$ is tested with several prize pairs, a common weight must
belong to all the resulting intervals. Empty intersection rejects that
combination of calibration and preference restrictions. Unknown utility
curvature and unknown weighting are not separately identified by simply
applying this formula with a convenient guessed utility.

\section*{What remains unresolved}
The results identify the exact ambiguity created by bounded mirror mismatch
and give a finite-panel calibration guarantee. They do not establish that
a particular mirror satisfies (1), identify arbitrary utility and weighting
jointly, or transport laboratory preferences to field choices with different
stakes and institutions. Those are central parts of the original agenda.
The bounds make the missing restriction explicit and permit reporting how
conclusions change as its tolerance $\eta$ varies.

\section{Completion Estimate}
45\%

This is a post hoc, heuristic estimate for the full catalogue target.
The attempt establishes observational nonidentification with systematic errors and derives sharp certainty-equivalent bounds and finite-panel guarantees under calibrated repeated mirrors. It still must validate mirror calibration, separate unknown utility from probability weighting, and establish transport and predictive performance for consequential field decisions.

\begin{thebibliography}{9}
\bibitem{oprea} R. Oprea (2024), \emph{Decisions under Risk Are Decisions
under Complexity}, American Economic Review 114(12),3789--3811.
\url{https://www.aeaweb.org/articles?id=10.1257/aer.20221227}.
\end{thebibliography}

\end{document}
